\documentclass[../DoAn.tex]{subfiles}
\begin{document}

Khóa luận đã giải quyết hai bài toán đặt ra ở Chương 1: (i) xây dựng một hệ thống web hoàn chỉnh hỗ trợ tạo, phân tích và tối ưu CV bằng mô hình ngôn ngữ lớn; (ii) thiết kế, hiện thực và đánh giá quy trình RAG hai bước phục vụ ứng viên tìm kiếm JD khớp yêu cầu. Chương này tóm tắt kết quả, nêu các hạn chế còn tồn tại, và đề xuất một số hướng phát triển trong tương lai.

\section{Tóm tắt kết quả đạt được}

\textbf{Về phần sản phẩm:}

\begin{enumerate}
\item Sinh viên đã xây dựng trình soạn thảo CV tương tác với layout hai cột (editor và preview), hỗ trợ chỉnh sửa theo từng phần, kéo thả sắp xếp thứ tự các mục, lưu tự động sau 2 giây, và xem trước định dạng A4 ngay khi đang sửa. Trình soạn thảo hỗ trợ xuất CV ra hai định dạng phổ biến: PDF và DOCX.

\item Sinh viên đã phát triển hệ thống web hoàn chỉnh với đăng nhập bằng email/mật khẩu, hỗ trợ phiên khách (không cần đăng nhập), quản lý nhiều CV trực tuyến cho từng người dùng, và giao diện song ngữ Việt--Anh. Hệ thống đã được kiểm thử qua 24 trường hợp kiểm thử chức năng và đều đạt yêu cầu.
\end{enumerate}

\textbf{Về phần thiết kế pipeline AI:}

\begin{enumerate}\setcounter{enumi}{2}
\item Sinh viên đã thiết kế và hiện thực luồng xử lý RAG hai bước phục vụ ứng viên tìm kiếm JD phù hợp. Pipeline gồm bước Hybrid Retrieval (TF-IDF + Embedding, lọc top-20 JD) và bước Structural + LLM Batch Ranking (xếp hạng top-5). Thiết kế hai bước cho phép mở rộng tới kho JD lớn mà không tăng chi phí gọi LLM tuyến tính theo kích thước kho. Phân tích lý do lựa chọn từng thành phần được trình bày trong Chương~5.

\item Sinh viên đã phân tích quyết định thiết kế then chốt \textbf{Hybrid Parser so với Parser thuần LLM} -- chỉ ra rằng cách dùng regex cho các trường có khuôn mẫu (liên hệ, kỹ năng) và LLM cho phần cần hiểu ngữ nghĩa (kinh nghiệm) là lựa chọn hợp lý về cả chất lượng đầu ra lẫn chi phí token, đồng thời giảm phụ thuộc vào LLM cho các trường đơn giản.
\end{enumerate}

\section{Hạn chế}

\textbf{Hạn chế về quy mô đánh giá.} Đánh giá quy trình RAG trong Chương~5 được thực hiện trên bộ dữ liệu Kaggle Resume Dataset với ground truth định nghĩa theo ngành nghề -- một proxy hợp lý nhưng chưa phải annotation thủ công. Khóa luận chưa có tập dữ liệu CV--JD tiếng Việt được gắn nhãn thủ công để xác nhận chất lượng pipeline trong ngữ cảnh tuyển dụng Việt Nam.

\textbf{Hạn chế của quy trình RAG.} Embedding tổng quát (\texttt{text-embedding-3-small}) chưa được fine-tune cho domain tuyển dụng. Structural Score hiện chỉ đánh giá ba tiêu chí cứng (kinh nghiệm, bằng cấp, kỹ năng); các yếu tố mềm phụ thuộc hoàn toàn vào LLM. Chất lượng kết quả cũng phụ thuộc vào chất lượng kho JD được crawl.

\textbf{Hạn chế về phạm vi ngôn ngữ.} Hệ thống hỗ trợ tiếng Việt và tiếng Anh nhưng các mô hình LLM và embedding được tối ưu chủ yếu cho tiếng Anh, nên chất lượng phân tích CV tiếng Việt có thể thấp hơn trong một số tình huống phức tạp.

\textbf{Hạn chế về tính năng hệ thống.} Hệ thống hiện chỉ xuất CV ra PDF và DOCX; chưa hỗ trợ các định dạng khác như LaTeX hay TXT. Module phân tích CV chưa hỗ trợ tệp scan PDF (cần OCR).

\section{Hướng phát triển trong tương lai}

\textbf{Mở rộng đánh giá sang tiếng Việt.} Thí nghiệm trong Chương~5 được thực hiện trên 229 CV và 307 JD tiếng Anh (14 ngành) từ Kaggle Resume Dataset \cite{kaggle_resume2023}. Hướng phát triển tiếp theo là thu thập tập dữ liệu CV--JD tiếng Việt được gắn nhãn thủ công để xác nhận chất lượng pipeline trong ngữ cảnh tuyển dụng Việt Nam, nơi các mô hình embedding và LLM hiện tại chưa được tối ưu.

\textbf{Cải thiện hỗ trợ tiếng Việt.} Thu thập thêm CV--JD tiếng Việt để kiểm thử và cải thiện chất lượng phân tích, đặc biệt với các mô hình embedding và LLM chưa được tối ưu cho tiếng Việt.

\textbf{Tích hợp OCR cho CV scan.} Tích hợp một mô hình OCR (như Tesseract hoặc PaddleOCR) để xử lý các CV đã được scan thành ảnh -- một dạng phổ biến trong thực tế tuyển dụng.

\textbf{Cải thiện gợi ý chỉnh sửa CV.} Phát triển một module riêng dùng LLM để gợi ý cách viết lại các mục trong CV cho phù hợp với một mô tả công việc cụ thể, kết hợp kết quả của Embedding (chấm điểm tổng thể) và LLM (gợi ý sửa chi tiết).

\textbf{Fine-tune embedding cho domain tuyển dụng tiếng Việt.} Embedding tổng quát hiện tại không được tối ưu cho ngữ cảnh tuyển dụng, đặc biệt với thuật ngữ kỹ thuật tiếng Việt. Hướng phát triển là thu thập tập dữ liệu cặp (CV, JD) có nhãn lớn hơn và fine-tune mô hình embedding với contrastive learning, tương tự cách các mô hình như BGE-M3 được huấn luyện cho tiếng Trung.

\textbf{Mở rộng luồng xử lý RAG tìm kiếm JD.} Bổ sung tính năng lọc cứng trước bước retrieval (theo địa điểm, mức lương, loại hợp đồng) để giảm kích thước corpus và tăng precision. Tích hợp feedback vòng lặp: ứng viên đánh giá kết quả (ứng tuyển/bỏ qua) để cải thiện xếp hạng JD theo thời gian.

\textbf{Phát triển sản phẩm thương mại.} Hệ thống hiện tại đã đủ tính năng cốt lõi để mang ra phục vụ người dùng thật. Hướng tiếp theo gồm: thêm gói trả phí (xuất nhiều CV/tháng, tính năng AI nâng cao), tích hợp thanh toán, và triển khai trên hạ tầng đám mây có khả năng mở rộng tự động.

\end{document}
